\chapter{Vacuum Energy and de Sitter Horizons}

\lectureoverview{Vacuum energy is unlike ordinary matter: it neither selects a preferred inertial frame nor dilutes as the universe expands. Once gravity is included, that apparently innocuous zero of energy changes the geometry of spacetime. We shall derive the matter-dominated law \(a(t)\propto t^{2/3}\), contrast it with the vacuum-dominated law \(a(t)\propto e^{Ht}\), and then rewrite de Sitter spacetime in conformal coordinates. The resulting causal diagram reveals an observer-dependent event horizon at fixed proper distance.}

\section{Quantum Fluctuations in the Vacuum}

We begin with the question that motivates the entire lecture. Quantum field theory assigns energy to empty space. Does this mean that real electron--positron pairs are continually emerging from nothing?

\begin{classroomqa}
\audiencequestion{Do electron--positron production and annihilation actually occur in the vacuum, or is that merely a bookkeeping device?\lecturetimestamp{00:00:06}}
\lecturerresponse{Virtual pairs contribute to the vacuum energy, but they are not ordinary particles that can be watched emerging from empty space. They are fluctuations of the quantum ground state and persist only within the limits appropriate to a virtual process.}
\end{classroomqa}

The harmonic oscillator gives the cleanest model. Classically it can remain motionless at the minimum of its potential. Quantum mechanically, exact position at the minimum and exact zero momentum are incompatible. The ground state compromises between the two and retains the zero-point energy
\begin{equation}
E_0=\frac{1}{2}\hbar\omega.
\label{eq:l08-zero-point}
\end{equation}
This is the least energy the oscillator can have. It is not an excited state and does not represent a stream of observable particles.

A free field decomposes into normal modes, each behaving like an oscillator. The electromagnetic field in a cavity is a familiar example. Every mode contributes a term of the form in Equation~\eqref{eq:l08-zero-point}; electron, quark, and other quantum fields contribute their own ground-state fluctuations. Adding these contributions produces what we call the vacuum energy. The notorious mismatch between naive microscopic estimates and the observed cosmological scale is real, but it is not the subject of this lecture.

\begin{classroomqa}
\audiencequestion{Will the discrepancy of roughly \(120\) orders of magnitude be discussed again?\lecturetimestamp{00:02:59}}
\lecturerresponse{Not here. The present goal is to understand cosmic horizons: how the cosmological equations produce accelerated expansion and why accelerated expansion creates event horizons.}
\end{classroomqa}

\begin{classroomqa}
\audiencequestion{Is the ultraviolet--infrared relation from the previous lecture part of tonight's discussion?\lecturetimestamp{00:03:59}}
\lecturerresponse{It may be revisited later, but it is not part of the planned argument. The immediate path runs from vacuum energy to cosmology and then to horizons.}
\end{classroomqa}

\section{Why Vacuum Energy Is Lorentz Invariant}

The next question appears elementary but exposes the defining property of vacuum energy. If empty space has a positive energy density, why does it not behave like an optical medium and alter the speed of light?

\begin{classroomqa}
\audiencequestion{Should a positive vacuum-energy density have a refractive effect, making light travel slightly more slowly?\lecturetimestamp{00:04:35}}
\lecturerresponse{Ordinary material can do that because it defines a preferred rest frame. Vacuum energy is different: it does not select such a frame, and every inertial observer obtains the same vacuum-energy density. Local measurements of light therefore retain the invariant speed \(c\).}
\end{classroomqa}

We must distinguish a symmetry of the laws from a symmetry of a particular state. The equations of physics may be translation invariant even though Stanford is not Berkeley and Earth is not interstellar space. Translating the \emph{entire} physical arrangement leaves the laws unchanged; the actual arrangement of matter need not look the same at every location. Likewise, rotationally invariant equations do not make this classroom rotationally symmetric.

The same distinction applies to Lorentz transformations. A gas has a rest frame. An observer moving through it sees moving molecules and a different energy density, so the state of the gas is not Lorentz invariant even though its governing equations are. Ordinary matter, radiation baths, and laboratory apparatus all break the symmetry by selecting special motions or directions.

Vacuum energy is the exceptional form of matter that does not do so. Imagine a detector that reports the vacuum-energy density. It gives some value \(\rho_{\rm vac}\) while at rest. Accelerate the detector and repeat the measurement: the result is unchanged. In relativistic notation this property is encoded by a stress tensor proportional to the metric,
\begin{equation}
T_{\mu\nu}^{\rm vac}=-\rho_{\rm vac}g_{\mu\nu},
\qquad
p_{\rm vac}=-\rho_{\rm vac}.
\label{eq:l08-vacuum-stress}
\end{equation}
\editorialnote{The pressure relation is not derived on the board. It is the covariant statement behind the lecture's two defining properties: Lorentz invariance and constant density during expansion.}

The local qualification matters. A light ray passing an observer is measured with nearby clocks and rulers to move at \(c\). Asking for the speed of a faraway light ray relative to us in an expanding geometry requires a choice of coordinates and a prescription for comparing distant events. That coordinate speed need not equal \(c\), even though every local observer measures \(c\).

\begin{classroomqa}
\audiencequestion{In this discussion, are vacuum energy and the cosmological constant synonymous?\lecturetimestamp{00:11:57}}
\lecturerresponse{Yes. They are two placements of the same term in Einstein's equations: it may be described as a cosmological constant on the geometric side or as Lorentz-invariant vacuum stress-energy on the source side.}
\end{classroomqa}

\section{What Dilutes and What Does Not}

Fill a large region with nonrelativistic particles of mass \(m\). On scales large enough to average over galaxies and clusters, suppose the distribution is homogeneous. If a comoving volume contains \(N\) particles, its matter density is
\begin{equation}
\rho_m=\frac{Nm}{V}.
\end{equation}
When every linear dimension doubles, the volume grows by \(2^3=8\), while \(N\) and \(m\) remain unchanged. Thus the density falls to one eighth of its former value. Ordinary matter dilutes because the same particles occupy a larger physical volume.

Radiation also dilutes, and more rapidly. Its number density falls with the volume while each photon's wavelength is stretched and its energy redshifted. We shall later record the result \(\rho_r\propto a^{-4}\), but we will not solve the radiation-dominated cosmology tonight.

Vacuum energy is the exception:
\begin{equation}
\rho_{\rm vac}(t)=\text{constant}.
\label{eq:l08-rho-vac-constant}
\end{equation}
It does not cluster, select a direction, select an inertial frame, or become less dense as space expands. If a comoving box grows, the vacuum energy contained in it grows in proportion to its volume.

\editorialnote{Constant density does not mean constant total vacuum energy inside an expanding comoving box. Covariant stress-energy conservation,
\(\dot\rho+3H(\rho+p)=0\), is satisfied because vacuum energy has negative pressure \(p_{\rm vac}=-\rho_{\rm vac}\). The work term supplies precisely the change in the box's total vacuum energy.}

\section{Gravity Sees the Zero of Energy}

In nongravitational laboratory physics, shifting every energy by the same constant usually changes nothing. Spectra depend on differences: the energy between an excited and an unexcited atom, or between a state with and without a particle. Empty space seems to offer no reference against which its absolute energy could be measured.

Gravity changes the question. Mass is energy, energy sources gravity, and therefore vacuum energy gravitates. A constant that was invisible in an atomic transition can curve spacetime.

Einstein's equation separates geometry from sources:
\begin{equation}
G_{\mu\nu}=8\pi G_N T_{\mu\nu},
\label{eq:l08-einstein}
\end{equation}
in units with \(c=1\). The Einstein tensor \(G_{\mu\nu}\) is built from the metric \(g_{\mu\nu}(x)\) and its derivatives. The stress tensor \(T_{\mu\nu}\) contains energy density, momentum density, pressure, and stress from matter and nongravitational fields. Vacuum energy belongs on this source side just as ordinary matter does.\editorialnote{The spoken introduction briefly associates \(8\pi/3\) with the tensor equation. The factor \(1/3\) belongs to the spatially flat Friedmann equation below, not to Equation~\eqref{eq:l08-einstein}.}

\begin{figure}[htbp]
\centering
\includegraphics[width=0.62\textwidth]{lecture_08_figure_02.png}
\caption{The geometric side of Einstein's equation, built from the spacetime metric \(g_{\mu\nu}(x)\).}
\lectureframe{00:21:58.820}
\end{figure}

\section{Geometry Reduced to a Scale Factor}

We now choose the simplest homogeneous expanding geometry with flat spatial slices:
\begin{equation}
ds^2=-dt^2+a(t)^2\,d\mathbf{x}^2,
\qquad
d\mathbf{x}^2=dx^2+dy^2+dz^2.
\label{eq:l08-frw-metric}
\end{equation}
The coordinate marks \(\mathbf{x}\) are painted on an infinite rubber sheet. They need not move, but their physical separations do. Two comoving points separated by \(\Delta\mathbf{x}\) have proper distance at fixed \(t\)
\begin{equation}
D(t)=a(t)\lvert\Delta\mathbf{x}\rvert.
\label{eq:l08-proper-distance}
\end{equation}
Thus all the time dependence of this highly symmetric geometry is stored in one function, the scale factor \(a(t)\). Constant \(a\) gives flat Minkowski spacetime after a spatial rescaling. Time-dependent \(a\) can leave each spatial slice flat while making the full spacetime curved.

\begin{classroomqa}
\audiencequestion{For coordinate marks \(x=7\) and \(x=13\), is their coordinate separation \(4\)?\lecturetimestamp{00:27:40}}
\lecturerresponse{No. It is \(13-7=6\), so their physical separation is \(6a(t)\). The correction is useful: \(x\) labels comoving position, while multiplication by \(a(t)\) gives physical distance.}
\end{classroomqa}

For this metric and a spatially uniform energy density, one component of Einstein's equation becomes the spatially flat Friedmann equation,
\begin{equation}
\left(\frac{\dot a}{a}\right)^2
=\frac{8\pi G_N}{3}\rho(t).
\label{eq:l08-friedmann}
\end{equation}
The ratio
\begin{equation}
H(t)\equiv\frac{\dot a(t)}{a(t)}
\label{eq:l08-hubble-definition}
\end{equation}
is the Hubble parameter. It is a definition at this stage and may vary with time.

\begin{figure}[htbp]
\centering
\includegraphics[width=0.90\textwidth]{lecture_08_figure_03.png}
\caption{The flat expanding metric and Friedmann equation after the missing factor \(1/a^2\) has been restored on the board.}
\lectureframe{00:30:30.000}
\end{figure}

\paragraph{The energy-balance picture.}
For this energy-balance picture, Equation~\eqref{eq:l08-friedmann} is written schematically as
\begin{equation}
-\frac{3}{8\pi G_N}H^2+\rho=0.
\label{eq:l08-energy-heuristic}
\end{equation}
The matter energy is positive, while the gravitational degree of freedom represented by the expanding scale factor carries a negative kinetic term. In this reduced problem, negative gravitational energy plus positive matter energy equals zero is a useful mnemonic.

\editorialnote{Equation~\eqref{eq:l08-energy-heuristic} is the Hamiltonian constraint in a highly symmetric minisuperspace model, not a coordinate-independent theorem that every arbitrary universe possesses a globally conserved total energy equal to zero. The lecture uses it as a heuristic interpretation of the Friedmann constraint.}

\section{Matter-Dominated Expansion}

Return to a comoving mathematical box with coordinate side \(\Delta x=1\). Its physical side is \(a(t)\), so its physical volume is \(a(t)^3\). If it contains \(N\) nonrelativistic particles of mass \(m\), then
\begin{equation}
\rho_m(t)=\frac{Nm}{a(t)^3}
=\frac{\rho_0}{a(t)^3}.
\label{eq:l08-matter-density}
\end{equation}
The particles themselves do not expand; their mean separation does.

\begin{figure}[htbp]
\centering
\includegraphics[width=0.84\textwidth]{lecture_08_figure_04.png}
\caption{The metric and Friedmann equation as the comoving-box argument is assembled.}
\lectureframe{00:36:07.940}
\end{figure}

Substitution into Equation~\eqref{eq:l08-friedmann} gives
\begin{equation}
\left(\frac{\dot a}{a}\right)^2
=\frac{8\pi G_N\rho_0}{3a^3}.
\end{equation}
Write \(K=8\pi G_N\rho_0/3\) and choose the expanding branch:
\begin{align}
\dot a^2&=\frac{K}{a},\\
\dot a&=\sqrt{K}\,a^{-1/2},\\
a^{1/2}\,da&=\sqrt{K}\,dt.
\label{eq:l08-matter-separated}
\end{align}

\begin{classroomqa}
\audiencequestion{How should the separated equation be solved?\lecturetimestamp{00:38:15}}
\lecturerresponse{Integrate both sides. The integral of \(a^{1/2}\,da\) is \(\frac{2}{3}a^{3/2}\), not \(\frac{3}{2}a^{3/2}\); \(\frac{3}{2}\) is the new power, while \(\frac{2}{3}\) is the coefficient.}
\end{classroomqa}

Accordingly,
\begin{equation}
\frac{2}{3}a^{3/2}
=\sqrt{K}\,(t-t_B),
\end{equation}
where \(t_B\) is an integration constant fixing the chosen origin of time. Apart from normalization and the time shift,
\begin{equation}
a(t)\propto (t-t_B)^{2/3}.
\label{eq:l08-matter-solution}
\end{equation}

\begin{figure}[htbp]
\centering
\includegraphics[width=0.86\textwidth]{lecture_08_figure_05.png}
\caption{The matter-density substitution and separation of variables leading to \(a^{1/2}da\propto dt\).}
\lectureframe{00:38:30.000}
\end{figure}

The curve \(t^{2/3}\) grows more slowly than a straight line. Its slope decreases, so a matter-dominated universe decelerates. In the Newtonian language used in the lecture, the particles began by separating, and their mutual gravitational attraction continually reduces the rate of separation.

\begin{figure}[htbp]
\centering
\includegraphics[width=0.86\textwidth]{lecture_08_figure_06.png}
\caption{The board comparison of linear expansion with the concave-down matter law \(a(t)\propto t^{2/3}\).}
\lectureframe{00:40:00.000}
\end{figure}

For orientation, a radiation bath has
\begin{equation}
\rho_r\propto a^{-4},
\end{equation}
rather than \(a^{-3}\): the volume dilution contributes three powers and cosmological redshift contributes one more. The corresponding scale factor can also be solved, but the lecture deliberately leaves that exercise aside.

\section{Vacuum-Dominated Expansion}

Now put only vacuum energy on the right-hand side of the Friedmann equation. Since \(\rho_{\rm vac}\) is constant,
\begin{equation}
\left(\frac{\dot a}{a}\right)^2
=\frac{8\pi G_N}{3}\rho_{\rm vac}
\equiv H^2.
\label{eq:l08-de-sitter-friedmann}
\end{equation}
There are two uses of \(H\) to keep separate. Equation~\eqref{eq:l08-hubble-definition} defines the possibly time-dependent quantity \(\dot a/a\). Equation~\eqref{eq:l08-de-sitter-friedmann} says that in a pure vacuum-energy universe its value is the constant
\begin{equation}
H=\sqrt{\frac{8\pi G_N\rho_{\rm vac}}{3}}.
\end{equation}
Taking the expanding branch yields
\begin{equation}
\dot a=Ha,
\end{equation}
and therefore
\begin{equation}
a(t)=C e^{Ht}.
\label{eq:l08-exponential-solution}
\end{equation}
The integration constant \(C\) fixes the arbitrary normalization of comoving coordinates. Growth proportional to the current value of \(a\) is exponential growth.

\begin{figure}[htbp]
\centering
\includegraphics[width=0.88\textwidth]{lecture_08_figure_07.png}
\caption{The vacuum-dominated equation \((\dot a/a)^2=H^2\), its first-order form \(\dot a=Ha\), and the exponential solution.}
\lectureframe{00:47:45.000}
\end{figure}

The Hubble law for comoving objects is
\begin{equation}
v_{\rm rec}=HD.
\label{eq:l08-hubble-law}
\end{equation}
In this pure de Sitter case, the proper distance at which the recession rate equals \(c\) is constant:
\begin{equation}
D_H=\frac{c}{H},
\label{eq:l08-hubble-radius}
\end{equation}
or \(D_H=H^{-1}\) in the \(c=1\) units used on the board. This fixed scale will reappear geometrically as the de Sitter event-horizon radius.\editorialnote{A Hubble radius is not automatically an event horizon in an arbitrary time-dependent cosmology. The identification works here because \(H\) is constant in exact de Sitter space.}

The lecture offers a thought experiment for a vacuum-energy detector. Place two very light test bodies far from other sources, initially at rest relative to the comoving flow, and make their mutual attraction negligible. Their separation obeys the exponential expansion. The effect is far too small for a terrestrial apparatus, but cosmic distances and lookback time make the same physics observable.

\section{From Matter Domination to Vacuum Domination}

The real cosmological model considered in the lecture contains both terms:
\begin{equation}
H(t)^2
=\frac{8\pi G_N}{3}
\left(\frac{\rho_0}{a(t)^3}+\rho_{\rm vac}\right).
\label{eq:l08-mixed-friedmann}
\end{equation}
At small \(a\), the matter term dominates because \(a^{-3}\) is large, and the solution resembles \(t^{2/3}\). As the universe grows, matter dilutes while vacuum energy remains fixed. Eventually the constant term wins and the scale factor bends upward toward \(e^{Ht}\).

\begin{figure}[htbp]
\centering
\includegraphics[width=0.87\textwidth]{lecture_08_figure_08.png}
\caption{The qualitative transition drawn in the lecture: early matter-dominated deceleration followed by late vacuum-dominated exponential growth.}
\lectureframe{00:53:30.000}
\end{figure}

\begin{classroomqa}
\audiencequestion{The exponential solution contains an integration constant. How do we know that it is not extremely small?\lecturetimestamp{00:52:41}}
\lecturerresponse{We measure it. Looking farther into the universe means looking farther into the past, so observations at different distances reconstruct the earlier portion of the expansion curve rather than merely its value today.}
\end{classroomqa}

The lecture reports that the then-available late-time observations were consistent with an exponential branch at roughly the five-percent level and describes the evidence as convincing, though not as exact as a precision measurement such as the electron's properties. This should be read as the lecture's historical assessment rather than as a present-day precision statement.

\begin{classroomqa}
\audiencequestion{Does the upward-bending expansion curve eventually prevent us from tracing observations back to the Big Bang or the cosmic microwave background?\lecturetimestamp{00:55:05}}
\lecturerresponse{That is not pursued as a question about extrapolating the graph. The relevant invariant issue is the existence of an event horizon, and that is the causal question we now turn to.}
\end{classroomqa}

\section{Bound Systems and the Distant Future}

If exponential expansion continues, every component that dilutes becomes negligible compared with \(\rho_{\rm vac}\). This does not mean that atoms, solar systems, or galaxies themselves are stretched apart. Systems held together strongly enough to remain bound do not follow the Hubble flow.

\begin{classroomqa}
\audiencequestion{Do galaxies that are not expanding today eventually expand when vacuum energy dominates?\lecturetimestamp{00:57:44}}
\lecturerresponse{Bound systems remain bound. The Milky Way and Andromeda belong to the nearby structure that will remain together, while sufficiently distant galaxies already moving with the expansion continue outward and eventually pass beyond our horizon.}
\end{classroomqa}

The distinction between \emph{attracted} and \emph{bound} is essential. A rocket launched above escape velocity remains under Earth's attractive gravity, yet it does not return. Likewise, ordinary gravitational attraction may be the largest local force between two galaxies without their relative orbit being bound.

\begin{classroomqa}
\audiencequestion{Our galaxy is gravitationally associated with a group. Does accelerated expansion eventually overcome that attraction?\lecturetimestamp{00:59:17}}
\lecturerresponse{Only the genuinely bound part of the local structure stays together. A cluster as a descriptive label is not enough; binding means that the relative motion lies below escape. More distant galaxies can be attracted and nevertheless recede above escape speed.}
\end{classroomqa}

For intuition, the lecture represents the relative radial acceleration schematically as an attractive inverse-square term plus a vacuum term linear in separation:
\begin{equation}
\ddot r\sim-\frac{G_NM}{r^2}+H^2r.
\label{eq:l08-newtonian-vacuum}
\end{equation}
The coefficients are time independent in the idealized de Sitter limit. The second term grows only because \(r\) grows; it does not arise from a gradual weakening of \(G_N\).

\editorialnote{Equation~\eqref{eq:l08-newtonian-vacuum} is a Newtonian relative-acceleration heuristic, not an additional force law fundamental to general relativity. In exact de Sitter space, neighboring freely falling worldlines have relative acceleration proportional to \(H^2r\).}

\begin{figure}[htbp]
\centering
\includegraphics[width=0.88\textwidth]{lecture_08_figure_09.png}
\caption{The board comparison between inverse-square attraction and the effective vacuum repulsion that grows linearly with separation.}
\lectureframe{01:04:10.000}
\end{figure}

\begin{classroomqa}
\audiencequestion{Does the outward effect itself increase with time?\lecturetimestamp{01:01:18}}
\lecturerresponse{Not at a fixed separation. Its coefficient is constant. Freely receding objects move farther apart, and the term proportional to \(r\) is then larger because their distance is larger.}
\end{classroomqa}

\begin{classroomqa}
\audiencequestion{Can the net effect be described as a gradual decrease of Newton's constant?\lecturetimestamp{01:02:30}}
\lecturerresponse{No. \(G_N\) remains constant. At a specified separation both coefficients are unchanged; only the relative importance of the \(1/r^2\) attraction and the term proportional to \(r\) changes with distance.}
\end{classroomqa}

\begin{classroomqa}
\audiencequestion{Were the same galaxies gravitationally bound in the early universe before vacuum energy became important?\lecturetimestamp{01:05:03}}
\lecturerresponse{Dominance of ordinary attraction does not by itself imply binding. Many galaxies were already separating at about escape speed or above it. Vacuum energy later determines whether the large-scale expansion decelerates toward recollapse or accelerates away, but largest force and bound orbit remain different statements.}
\end{classroomqa}

Thus a far-future observer may inhabit a merged local structure while every unbound galaxy has become unobservable. Without historical records, that observer could reasonably infer an almost empty universe. This causal loneliness, rather than destruction of local objects, is the physical lead-in to the horizon.

\section{What an Event Horizon Means}

After the break, set
\begin{equation}
a(t)=e^{Ht}
\end{equation}
and write the metric as
\begin{equation}
ds^2=-dt^2+e^{2Ht}d\mathbf{x}^2.
\label{eq:l08-de-sitter-metric}
\end{equation}

\begin{classroomqa}
\audiencequestion{What is the name of the spacetime described by this metric?\lecturetimestamp{01:08:13}}
\lecturerresponse{It is de Sitter spacetime, commonly shortened to de Sitter space. Its horizons matter because a vacuum-dominated universe approaches this geometry at late times.}
\end{classroomqa}

Before drawing its causal structure, recall the role of a Penrose diagram. Infinite intervals are compressed into a finite picture while null rays remain at \(45^\circ\). In flat spacetime, massive objects end at future timelike infinity, null rays end at null infinity, and an inertial observer allowed to wait arbitrarily long has no event horizon excluding a permanent region of spacetime.

\begin{figure}[htbp]
\centering
\includegraphics[width=0.80\textwidth]{lecture_08_figure_10.png}
\caption{Flat-spacetime Penrose diagram used to distinguish timelike and null destinations before introducing a horizon.}
\lectureframe{01:10:35.000}
\end{figure}

A black-hole diagram supplies the contrast. An exterior observer can receive signals only from events in the causal past of the observer's future worldline. Even as the observation time tends to infinity, the interior region remains excluded. The boundary between events that can communicate with that observer and events that never can is the event horizon.

\begin{definition}
The future event horizon of an observer is the boundary of the set of events from which a future-directed causal signal can ever reach that observer's worldline.
\label{def:l08-event-horizon}
\end{definition}

\begin{figure}[htbp]
\centering
\includegraphics[width=0.80\textwidth]{lecture_08_figure_11.png}
\caption{The black-hole comparison: backward light cones from arbitrarily late exterior observations never include the region beyond the event horizon.}
\lectureframe{01:13:50.000}
\end{figure}

An event horizon is therefore a global causal boundary, not merely the distance at which an object happens to look faint today. To locate it, one must know which signals can arrive even after an arbitrarily long wait.

\section{Compressing the Infinite Future}

The spatial coordinates in Equation~\eqref{eq:l08-de-sitter-metric} already range over an infinite flat slice. We shall leave them alone and compress only time. Introduce conformal time \(T\) by requiring
\begin{equation}
dt^2=e^{2Ht}dT^2.
\label{eq:l08-conformal-definition}
\end{equation}
Taking the future-directed square root and rearranging gives
\begin{equation}
dT=e^{-Ht}dt.
\end{equation}

\begin{classroomqa}
\audiencequestion{What is the integral, including its sign?\lecturetimestamp{01:17:25}}
\lecturerresponse{It is \(-H^{-1}e^{-Ht}\). An additive constant merely shifts the origin of \(T\), so choose it to place the infinite future at \(T=0\).}
\end{classroomqa}

Thus
\begin{equation}
T=-\frac{1}{H}e^{-Ht},
\qquad -\infty<T<0.
\label{eq:l08-conformal-time}
\end{equation}
As \(t\to-\infty\), \(T\to-\infty\). As the proper time \(t\to+\infty\), \(T\to0^{-}\). An infinite future has been placed on the finite horizontal line \(T=0\); no observer reaches it in finite proper time.

\begin{figure}[htbp]
\centering
\includegraphics[width=0.86\textwidth]{lecture_08_figure_12.png}
\caption{The conformal-time integral and the mapping of the infinite future \(t\to+\infty\) to \(T=0^{-}\).}
\lectureframe{01:19:10.000}
\end{figure}

From Equation~\eqref{eq:l08-conformal-time},
\begin{equation}
HT=-e^{-Ht},
\qquad
e^{2Ht}=\frac{1}{H^2T^2}.
\label{eq:l08-conformal-factor}
\end{equation}
Using \(dt=e^{Ht}dT\) in Equation~\eqref{eq:l08-de-sitter-metric} then yields
\begin{equation}
ds^2=\frac{1}{H^2T^2}
\left(-dT^2+d\mathbf{x}^2\right).
\label{eq:l08-conformal-metric}
\end{equation}
The expression in parentheses is the Minkowski metric. All the de Sitter geometry now sits in a common conformal factor that diverges as \(T\to0^{-}\).

\begin{figure}[htbp]
\centering
\includegraphics[width=0.87\textwidth]{lecture_08_figure_13.png}
\caption{The algebra relating \(e^{2Ht}\) to \(1/(H^2T^2)\), preparing the conformal form of the metric.}
\lectureframe{01:22:50.000}
\end{figure}

\begin{classroomqa}
\audiencequestion{What condition determines the path of a light ray in this metric?\lecturetimestamp{01:23:49}}
\lecturerresponse{A light ray follows a null path, \(ds^2=0\). The common conformal factor therefore cancels, leaving \(dT=\pm dx\) along a chosen spatial direction.}
\end{classroomqa}

Light rays are straight \(45^\circ\) lines in the \((T,x)\) diagram:
\begin{equation}
\frac{dx}{dT}=\pm1.
\label{eq:l08-null-lines}
\end{equation}
The coordinate picture is highly compressed near \(T=0\), but it preserves precisely what we need: which events can communicate.

\begin{figure}[htbp]
\centering
\includegraphics[width=0.86\textwidth]{lecture_08_figure_14.png}
\caption{The conformal metric at the moment the null condition \(ds^2=0\) is used to recover \(45^\circ\) light rays.}
\lectureframe{01:24:05.000}
\end{figure}

\section{Alice and Bob in de Sitter Space}

Place Alice at fixed comoving position \(x=0\). Her worldline rises vertically and ends on the diagram at \(T=0\), although her wristwatch requires infinite proper time to get there. Draw the past-directed null rays from that endpoint. They bound the largest region from which Alice can ever receive a signal. Events outside it remain forever inaccessible.

\begin{figure}[htbp]
\centering
\includegraphics[width=0.88\textwidth]{lecture_08_figure_15.png}
\caption{Alice's limiting past light cone. Its null boundary is her observer-dependent de Sitter event horizon.}
\lectureframe{01:28:50.000}
\end{figure}

This horizon is private. A different observer has a different future endpoint and therefore a different limiting past light cone. Suppose Bob once lies within Alice's causal region and the two exchange messages. If Bob later follows the Hubble flow across Alice's horizon, signals emitted after that crossing cannot reach Alice, however long she waits. Bob has his own horizon and eventually loses the corresponding ability to receive late signals from Alice.

\begin{figure}[htbp]
\centering
\includegraphics[width=0.88\textwidth]{lecture_08_figure_16.png}
\caption{Alice and Bob with their distinct causal boundaries. A comoving Bob can leave Alice's causal past even though they exchanged signals earlier.}
\lectureframe{01:30:20.000}
\end{figure}

If Alice and Bob hold hands, or more realistically remain bound at fixed proper separation, their comoving-coordinate separation shrinks as \(a^{-1}\). Their curves approach each other in the conformal diagram, but this does not mean their proper distance vanishes. The growing conformal factor compensates for the shrinking coordinate gap, and communication remains possible indefinitely.

\begin{classroomqa}
\audiencequestion{Can Bob see Alice at the same time after their horizons separate them?\lecturetimestamp{01:31:43}}
\lecturerresponse{No observer ever sees a distant event at the same time in that literal sense; observation is always of an event in the past light cone. The sharper statement is that after horizon crossing, no newly emitted signal from the other worldline can arrive.}
\end{classroomqa}

In the far-future cosmological application, Alice represents our bound local remnant and Bob represents each unbound distant galaxy. The galaxies do not suddenly wink out at a finite instant on Alice's clock. Their later signals are stretched over longer intervals and shifted to longer wavelengths until they become physically undetectable.

\begin{classroomqa}
\audiencequestion{Does the theory really say that no distant galaxies will remain visible in the far future?\lecturetimestamp{01:33:38}}
\lecturerresponse{Unbound galaxies pass beyond the local observer's event horizon. Their final visible signals become increasingly delayed and redshifted, leaving only the bound local structure observable for practical purposes.}
\end{classroomqa}

\section{Redshift at the Horizon}

One can still draw an old light signal from Bob reaching Alice at a very late conformal time, which can make it appear that she continues to see him. The physical content lies in the spacing and energy of those signals. Equal intervals of Bob's clock arrive separated by ever larger intervals of Alice's proper time. Frequencies and photon energies fall:
\begin{equation}
\nu_{\rm received}\longrightarrow0,
\qquad
E_{\rm received}=h\nu_{\rm received}\longrightarrow0.
\end{equation}
Bob's image slows, reddens, and fades. The limiting signal carries vanishing usable information.

\begin{figure}[htbp]
\centering
\includegraphics[width=0.89\textwidth]{lecture_08_figure_17.png}
\caption{The late board state comparing de Sitter horizon redshift with the familiar external view of a black-hole horizon.}
\lectureframe{01:36:20.000}
\end{figure}

\begin{classroomqa}
\audiencequestion{Is this like a black hole, where an outside observer never sees an infalling object cross the horizon?\lecturetimestamp{01:36:09}}
\lecturerresponse{Yes. In both cases the observer receives increasingly redshifted signals and never directly sees the crossing completed. The causal diagram says which later emissions can never arrive; the redshift explains how the last observable information fades.}
\end{classroomqa}

The analogy prompts a question about temperature.

\begin{classroomqa}
\audiencequestion{How could Alice measure the temperature just above a horizon?\lecturetimestamp{01:39:34}}
\lecturerresponse{Operationally, she could lower a thermometer on a cable, keep it from crossing, and retrieve it. A de Sitter observer could imagine the analogous experiment near the cosmological horizon. Horizons are hot in the quantum theory, but the present lecture stops at the classical geometry and postpones that calculation.}
\end{classroomqa}

No formula for de Sitter temperature is derived here. Hawking radiation enters only as a preview of the quantum correction to this classical picture.

\section{The Fixed Horizon Radius}

The conformal diagram can be translated back into the original time coordinate. Two comoving points at fixed coordinate separation satisfy
\begin{equation}
D(t)=e^{Ht}\lvert\Delta x\rvert,
\end{equation}
so their proper distance grows exponentially.

\begin{classroomqa}
\audiencequestion{What does the same motion look like in the original, untransformed time coordinate?\lecturetimestamp{01:41:14}}
\lecturerresponse{For fixed comoving coordinates, the entire spatial metric grows with \(e^{2Ht}\); hence proper separations grow like \(e^{Ht}\). Undoing the coordinate transformation makes that expansion explicit.}
\end{classroomqa}

The horizon behaves differently because its comoving coordinate distance from Alice is not fixed. From the \(45^\circ\) null line ending at \(T=0\), the horizon distance on a slice of constant \(T<0\) is
\begin{equation}
\Delta x_H=-T=\frac{1}{H}e^{-Ht}.
\label{eq:l08-coordinate-horizon}
\end{equation}
Multiplying by the scale factor gives
\begin{equation}
D_H(t)=a(t)\Delta x_H
=e^{Ht}\frac{e^{-Ht}}{H}
=\frac{1}{H}.
\label{eq:l08-fixed-horizon}
\end{equation}
Restoring units, \(D_H=c/H\). The coordinate gap shrinks, the metric scale grows, and their product remains constant.

\begin{classroomqa}
\audiencequestion{As time passes, does the distance from Alice to the horizon grow, shrink, or stay the same?\lecturetimestamp{01:42:42}}
\lecturerresponse{It stays the same. Its coordinate distance shrinks while the conformal factor grows, leaving the proper horizon radius fixed at \(H^{-1}\) in units with \(c=1\).}
\end{classroomqa}

\begin{figure}[H]
\centering
\includegraphics[width=0.90\textwidth]{lecture_08_figure_18.png}
\caption{The final board state: successive horizontal slices have shrinking coordinate distance to the horizon but the same proper distance.}
\lectureframe{01:43:20.000}
\end{figure}

\begin{classroomqa}
\audiencequestion{Does Alice continue to see the past projected onto the horizon because the old images never quite disappear?\lecturetimestamp{01:43:45}}
\lecturerresponse{Classically, the images persist only in the limiting sense that they become harder and harder to detect through unbounded redshift. That statement omits the Hawking radiation associated with the horizon, which will modify the quantum description.}
\end{classroomqa}

The lecture therefore closes at a precise boundary. Vacuum energy has produced exponential expansion; exponential expansion has produced a fixed, observer-dependent event horizon; and signals near that horizon have been driven to vanishing frequency. The horizon's temperature and Hawking radiation belong to what comes next. Assuming, as Susskind says, that we do not cross one another's horizons, we meet again in the next lecture.
